\documentclass{article}
\usepackage[utf8]{inputenc}
\usepackage{natbib}
\usepackage{graphicx}

\title{Complex QA and Language Models Hybrid Architectures Survey}
\author{Your Name}
\date{\today}

\begin{document}

\maketitle

\begin{abstract}
This paper reviews the state-of-the-art of language model architectures and strategies for complex question-answering (CQA) with a focus on hybridization. It discusses challenges such as domain adaptation, decomposition strategies, and the need for explainability in complex QA tasks. Current solutions and emerging research trends in hybrid LLM architectures, training strategies, and multimodal search are analyzed.
\end{abstract}

\tableofcontents
\newpage

\section{Introduction}
Recent advancements in large language models (LLMs), such as ChatGPT and GALACTICA, have demonstrated their potential in leveraging public data to tackle standard question-answering tasks. However, when confronted with complex questions that require nuanced understanding across multiple domains and cultural contexts, LLMs often face limitations in accuracy, fairness, and robustness. This paper explores the evolution of LLM architectures and strategies aimed at enhancing their performance in complex QA scenarios.

\section{Required Skills and Evaluation Techniques}
To effectively address complex QA challenges, researchers and practitioners require a diverse set of skills and robust evaluation techniques. Skills encompass not only technical expertise in natural language processing (NLP) and machine learning but also domain-specific knowledge in fields such as sociology, economics, and environmental science. Evaluation techniques must go beyond traditional accuracy metrics to include fairness, robustness against adversarial attacks, and sensitivity to multi-modal inputs.

\section{Challenges in Complex QA}
Complex QA poses several challenges that standard LLMs struggle to overcome. These include domain adaptation, decomposition and efficient multi-step QA, long-form and non-factoid QA, and the need for safety and multi-sensitivity data protection. Furthermore, issues like hallucinations, explainability, truthfulness, and temporal reasoning add layers of complexity that require innovative solutions.

\section{Current Solutions and Promising Research Trends}
Recent advancements in hybrid LLM architectures have shown promise in enhancing complex QA capabilities. These architectures integrate techniques such as neuro-symbolic reasoning, structured knowledge grounding, and multimodal search to improve accuracy and robustness. Training strategies, including active human reinforcement learning supervised with AI, and program synthesis are also explored for their potential in tackling complex QA challenges.

\section{Conclusion}
In conclusion, while LLMs have made significant strides in standard QA tasks, their application to complex questions requires sophisticated architectural patterns and strategic training approaches. Future research efforts should focus on enhancing LLM capabilities in domain adaptation, explainability, and handling sensitive data to meet the diverse needs of complex QA applications.

\bibliographystyle{plain}
\bibliography{prompt1_try4_35} 

\end{document}
